\chapter{Background: Why Memory Operations Are Reordered}
\label{ch:background}

This chapter establishes the microarchitectural facts on which the rest of the
paper depends.
None of the mechanisms described here was introduced in order to
weaken the memory model; each was introduced to solve a performance problem, and
the weakening of ordering guarantees is a side effect.
That the side effect is
unavoidable --- that the speed and the reordering are two descriptions of the same
property --- is the central observation of the chapter.

\section{The Memory Wall}
\label{sec:memwall}

A processor core executes arithmetic at a rate of one or more operations per
clock cycle.
Access to main memory takes on the order of $200$ cycles.
Sarangi
develops this asymmetry in \S11.1.1 of the prescribed text
\cite{sarangi2025basic}; the consequence is that a core that waited for every
memory access would spend the overwhelming majority of its cycles idle.

Two structures address this, and they address different halves of the problem.
The \emph{cache} shortens read latency by keeping copies of recently accessed
lines close to the core.
The \emph{store buffer} eliminates write latency by
allowing the core to retire a store without waiting for it to reach memory at all.

\begin{figure}[H]
\centering
\begin{tikzpicture}[node distance=9mm]
  \node[core] (c)                       {Core\\\scriptsize registers, $\sim$1 cycle};
  \node[buf]  (sb) [right=22mm of c]    {Store buffer\\\scriptsize pending writes};
  \node[mem]  (l1) [below=12mm of c]    {L1 cache\\\scriptsize $\sim$4 cycles};
  \node[mem]  (l2) [below=of l1]        {L2 / L3\\\scriptsize $\sim$12--40 cycles};
  \node[mem]  (mm) [below=of l2]        {Main memory\\\scriptsize $\sim$200 cycles};

  \draw[flow] (c) -- node[lbl,left]{loads} (l1);
  \draw[flow] (c) -- node[lbl,above]{stores} (sb);
  \draw[slow] (sb) |- node[lbl,pos=0.75,above]{drained later} (l1.east);
  \draw[flow] (l1) -- (l2);
  \draw[flow] (l2) -- (mm);
  \draw[flow] (sb.north) .. controls +(0,8mm) and +(0,8mm) .. (c.north)
        node[lbl,pos=0.5,above]{forwarding};
\end{tikzpicture}
\caption{The memory path of a single core. Loads traverse the cache hierarchy;
stores are placed in the store buffer and drained asynchronously. The forwarding
path returns a pending store's value to a subsequent load from the same core.}
\label{fig:hierarchy}
\end{figure}

\section{Caches}
\label{sec:caches}

Caches exploit temporal and spatial locality (\S11.1.3--\S11.1.6 of the text) to
serve most accesses without a memory transaction.
For the purposes of this paper
three properties matter.

\paragraph{Caches are private.} In a multicore processor, the L1 and typically
the L2 cache belong to a single core.
A line may therefore be resident in several
caches simultaneously, and the copies may disagree unless something is done about
it.
That something is coherence, and it is the subject of Chapter~3.

\paragraph{Cache misses are non-blocking.} A modern L1 permits multiple
outstanding misses, tracked in miss-status holding registers.
Consequently two
loads issued in program order may complete in either order, depending on which
one hits.

\paragraph{Cache line granularity.} Coherence operates on lines, typically 64
bytes, not on individual words.
Two variables that are logically unrelated but
share a line will interact through the coherence protocol.
This phenomenon,
false sharing, will matter in Chapter~6 when litmus test variables are laid out.

\section{The Store Buffer}
\label{sec:storebuffer}

The store buffer is the mechanism most directly responsible for the behaviour
described in \S\ref{sec:motivating}, and it deserves careful statement.

When a core executes a store, the value to be written is already in hand.
Unlike
a load, the store produces no result that subsequent instructions require.
There
is therefore no data dependence forcing the core to wait for the store to
complete.
The store buffer exploits this: the store is placed in a small FIFO
queue --- typically 40 to 60 entries --- and the core proceeds immediately.
A
separate mechanism drains the queue into the cache in the background.

The cost avoided is larger than it first appears.
Writing to a cache line
requires that the line be held in a state permitting modification.
If it is not,
the core must first issue a coherence request for exclusive ownership and wait for
every other core holding a copy to relinquish it.
On a large multicore that
round-trip may take hundreds of cycles.
Without a store buffer, the core stalls
for all of it.

\subsection{Store-to-Load Forwarding}

If a core stores to an address and subsequently loads from the same address, the
correct value may still be sitting in the store buffer rather than in the cache.
The load therefore checks the store buffer before the cache, and if a matching
pending store is found the value is returned directly.
This is
\emph{store-to-load forwarding}.

Forwarding is why single-threaded programs are unaffected by any of this.
A core
always observes its own writes in program order, whatever the rest of the system
sees.
The uniprocessor programming model is preserved exactly; only the
inter-core view is altered.

\subsection{The Unavoidable Consequence}

We can now state the point precisely.

\begin{quote}
The store buffer is fast \emph{because} a store held in it is not yet visible to
other cores.
The invisibility is not an implementation defect that better
engineering could remove; it is the same fact as the performance benefit, viewed
from the other side.
Any mechanism that made pending stores immediately visible
would, by definition, require the core to wait for them, which is the stall the
store buffer exists to avoid.
\end{quote}

\noindent Figure~\ref{fig:sbtest} traces Listing~\ref{lst:sb} through this
structure and shows how the outcome $\texttt{r0} = \texttt{r1} = 0$ arises without
any component behaving incorrectly.

\begin{figure}[H]
\centering
\begin{tikzpicture}
  \node[core] (c0) at (0,3)      {Core 0\\\scriptsize \code{x=1; r0=y;}};
  \node[core] (c1) at (8,3)      {Core 1\\\scriptsize \code{y=1; r1=x;}};
  \node[buf]  (b0) at (0,1.4)    {Store buffer\\\scriptsize holds \code{x=1}};
  \node[buf]  (b1) at (8,1.4)    {Store buffer\\\scriptsize holds \code{y=1}};
  \node[mem]  (m)  at (4,-0.4)   {Shared memory\\\scriptsize \code{x==0}, \code{y==0}};

  \draw[flow] (c0) -- (b0);
  \draw[flow] (c1) -- (b1);
  \draw[slow] (b0) -- node[lbl,left,pos=0.4]{not yet\\drained} (m.west);
  \draw[slow] (b1) -- node[lbl,right,pos=0.4]{not yet\\drained} (m.east);

  \draw[flow] (m.west)  .. controls +(-3.4,0) and +(-1.6,-1.4) .. (c0.south west)
        node[lbl,pos=0.45,left]{reads 0};
  \draw[flow] (m.east)  .. controls +(3.4,0) and +(1.6,-1.4) .. (c1.south east)
        node[lbl,pos=0.45,right]{reads 0};
\end{tikzpicture}
\caption{Execution of Listing~\ref{lst:sb}. Each store is buffered locally while
the corresponding load is served from shared memory, which still holds the
initial values. Both loads return zero. No component has malfunctioned.}
\label{fig:sbtest}
\end{figure}

\section{Out-of-Order Execution}
\label{sec:ooo}

Store buffering relaxes the ordering of a store with respect to \emph{later}
operations.
Out-of-order execution, treated in \S10.11.4 of the prescribed text,
relaxes considerably more.

An out-of-order core fetches instructions in program order, issues them for
execution as soon as their operands are available regardless of order, and retires
them in program order so that exceptions remain precise.
The load-store queue
tracks memory operations, resolves addresses as they become known, and permits
loads to be issued speculatively ahead of older stores whose addresses are not yet
computed.

Three consequences follow for the memory model:

\begin{enumerate}
\item \textbf{Loads may execute out of order with respect to one another}, since a
  load that hits in the cache completes before an older load that misses.
\item \textbf{Loads may execute ahead of older stores.} This is a deliberate
  optimisation; a store's address may not be known for many cycles, and blocking
  all subsequent loads until it is would be very costly.
\item \textbf{A squash mechanism already exists.} When a branch is mispredicted,
  the core discards all speculative state after the branch and restarts.
This
  machinery is general and, as Chapter~4 and the discussion of speculative SC will
  show, can be repurposed to recover from ordering violations.
\end{enumerate}

The third point is the one on which this paper's argument turns.
The expensive
component of an optimistic implementation of sequential consistency --- the
ability to undo work --- is not an addition to a modern core.
It is already there.

\section{The Compiler}
\label{sec:compiler}

Hardware is not the only source of reordering.
An optimising compiler reasons
about one thread at a time and freely reorders, coalesces, hoists and eliminates
memory accesses that no single-threaded observer could distinguish.
Register
allocation alone can remove a shared-variable read entirely, converting a spin
loop into an infinite loop.

The relevance to this paper is twofold.
First, litmus tests must be written so
that the compiler does not optimise away the very accesses being measured; the
techniques for this are set out in Chapter~6.
Second, any claim about the cost of
sequential consistency must account for the compiler, since hardware SC without
compiler SC provides no guarantee to the programmer at all.
Marino et al.\ address
precisely this \cite{marino2011sc}, and Chapter~11 returns to it.

\section{Summary}
\label{sec:bg-summary}

\begin{table}[H]
\centering
\small
\caption{Mechanisms that reorder memory operations, and their motivation.}
\label{tab:mechanisms}
\begin{tabular}{@{}lll@{}}
\toprule
\textbf{Mechanism} & \textbf{Introduced to} & \textbf{Reordering permitted} \\
\midrule
Store buffer            & hide store latency        & store $\rightarrow$ later load \\
Non-blocking cache      & overlap misses            & load $\rightarrow$ load \\
Out-of-order issue      & extract parallelism       & load $\rightarrow$ load, load $\rightarrow$ store \\
Write coalescing        & reduce bus traffic        & store $\rightarrow$ store \\
Compiler optimisation   & reduce instruction count  & all of the above \\
\bottomrule
\end{tabular}
\end{table}

Every entry in Table~\ref{tab:mechanisms} is a performance mechanism whose
ordering consequence is incidental. The architectural question --- which of these
reorderings a machine will admit to performing, and which it will hide from the
programmer --- is the subject of Chapter~4.
First, however, Chapter~3 must dispose
of the mechanism that is commonly, and mistakenly, expected to solve the problem.
